\scanpage{3}
On a donc une suite exacte
\[
(6)\qquad\diagram{\begin{tikzcd}[ampersand replacement=\&]
1\arrow[r]\&\pi_1(X)\arrow[r]\arrow[d,equal]\&\mathcal E_X\arrow[r]\&\mathcal G_X\arrow[r]\arrow[d]\&1\\
\&\operatorname{Aut}(\widetilde X/X)\&\&\Gamma\&
\end{tikzcd}}
\]
\deleted{La donnée d'un point}

Considérons un point $\xi$
\[
(7)\qquad\xi\in X(\overline{\mathbf Q})
\]
et soit $\mathcal G_X(\xi)\subset\mathcal G_X$ le sous-groupe
de $\mathcal G_X$ qui fixe $\xi$. (NB. Si $G_X$ est
\margin{$X$ est de t.f. $/\overline{\mathbf Q}$ et \unsure{sép.}}
fini, $\mathcal G_X$ est profini et $\mathcal G_X(\xi)$ est un
s-groupe ouvert\unsure{\ldots}). On trouve alors,
si on choisit pour $\widetilde X$ le rev. universel
\inserted{$\widetilde X(\xi)$} pointé au dessus de $\xi$, un
splitting partiel de (6)
\[
(8)\qquad\diagram{\begin{tikzcd}[ampersand replacement=\&]
\mathcal E_X\arrow[r]\&\mathcal G_X\arrow[r,dash]\arrow[d,phantom,"\cup" description]\&1\\
\&\mathcal G_X(\xi)\arrow[ul]\&
\end{tikzcd}}
\]
Dans le cas où $\widetilde X$ est quelconque, on
ne trouve qu'un hom. mod conjugaison
\unsure{aux} autom intérieurs provenant de
$\pi_1(X)$. \deleted{Si $X$ est de t.f. et}
\mspage{2}
